\begin{introduction}

La boiterie désigne une altération de la locomotion, le plus souvent causée par la douleur
d'une lésion du pied ou du membre. Elle se reconnaît cliniquement à une démarche irrégulière,
un dos voûté, une répartition inégale du poids entre les membres et une réticence à se
déplacer~; son évaluation de référence repose sur une notation locomotrice standardisée
\citep{sprecher1997lameness,chapinal2010measurement}. La boiterie est un phénomène clinique.
Les changements comportementaux mesurés par les capteurs n'en constituent qu'un reflet
indirect. Cette distinction guide l'ensemble du mémoire.

La boiterie constitue un problème majeur de bien-être animal et de performance en élevage
laitier.
Les études épidémiologiques rapportent fréquemment des prévalences de l'ordre de 20\,\% ou
davantage, avec des variations selon la population, la saison et la définition clinique
retenue \citep{whay2003assessment,cook2003prevalence}. Ses
conséquences comprennent la douleur, la baisse de production, une fertilité dégradée, des
soins supplémentaires et des réformes anticipées
\citep{green2002financial,bruijnis2010assessing}. Les estimations publiées chiffrent la
perte de production laitière à environ 360~kg par lactation chez les vaches cliniquement
boiteuses et le coût moyen d'un trouble du pied à 95~dollars américains par cas (18~\$~US
s'il est subclinique), soit de l'ordre de 75~\$~US par vache et par an à l'échelle
du troupeau. Au-delà du coût, la douleur associée est
reconnue comme l'une des principales sources de souffrance chronique chez la vache laitière
\citep{whay2003assessment}. Une identification plus précoce peut donc orienter l'examen du
pied et la prise en charge avant l'installation d'un épisode chronique
\citep{ahdb_mobility_score_sheet_48h_score2}.

L'élevage de précision propose de compléter l'observation humaine par une surveillance
continue au moyen de capteurs portés \citep{bewley2015precision}. Cette surveillance ne
résout toutefois pas, à elle seule, le problème de terrain: les accéléromètres et podomètres
mesurent des comportements, non la douleur ou la lésion. Une alerte produite à partir des pas,
de la posture ou de l'intensité de mouvement ne peut ainsi être assimilée à un diagnostic
clinique \citep{alsaaod2019automatic,oleary2020accelerometers}. L'enjeu devient donc de
transformer des séries comportementales indirectes en priorités d'observation fiables,
traçables et utilisables par l'éleveur.

Dans cette perspective, la manifestation la mieux étayée dans les données disponibles est une
réduction durable de l'activité locomotrice, qui se traduit par la diminution des pas, du
mouvement et des
transitions, accompagnée de changements posturaux. Ces signaux restent néanmoins non
spécifiques. Des situations différentes peuvent produire une forme semblable, notamment une
autre maladie, la chaleur ou certaines conditions d'élevage. Une difficulté supplémentaire est
que certaines situations douloureuses
pourraient s'accompagner d'agitation ou d'instabilité avant une hypoactivité, mais cette
trajectoire n'est pas démontrée dans le corpus principal. Une hausse brute de l'indice de
mouvement (\textit{Motion Index}) peut aussi refléter l'exercice, l'œstrus, une manipulation ou
un artefact de capteur.

La réponse proposée ne remplace pas le jugement clinique: elle ordonne les vérifications. HYPO
repère d'abord les trajectoires d'hypoactivité persistantes; INSTABILITÉ examine séparément
l'apport des profils d'agitation et leur coût en vérifications.

Le mémoire répond à une question principale: \textbf{comment détecter de façon reproductible
des dégradations comportementales graduelles et persistantes à partir de séries IceQube non
étiquetées sans les assimiler à un diagnostic de boiterie?} Il examine également une question
secondaire: \textbf{une branche d'instabilité peut-elle apporter une information complémentaire
sans dégrader excessivement la localisation temporelle ni augmenter la charge de vérification?}

L'objectif principal est de concevoir et d'évaluer un système reproductible qui transforme les
signaux IceQube en priorités de vérification comportementale. HYPO repère les baisses graduelles
et persistantes; INSTABILITÉ examine un apport complémentaire éventuel sans dégrader
excessivement la localisation temporelle ni augmenter la charge de vérification.

Pour atteindre cet objectif, la démarche méthodologique suit six étapes:
\begin{enumerate}[itemsep=0pt, parsep=0pt, topsep=0.3\baselineskip]
  \item construire une période de référence chronologique propre à chaque vache et à chaque créneau
  horaire;
  \item détecter les baisses multivariées persistantes par une branche HYPO;
  \item évaluer HYPO au niveau de l'événement et de la vache, sur des événements postérieurs à la
  période de référence, avec une ablation contre Isolation Forest, LOF et un comparateur pédométrique,
  et une analyse de sensibilité des paramètres;
  \item explorer une branche INSTABILITÉ, des persistances distinctes et cinq règles de
  fusion;
  \item mesurer explicitement les contrôles brefs, les facteurs confondants, la charge de
  fond et le coût d'exécution de la chaîne de traitement complète;
  \item confronter, sans réglage a posteriori, les sorties à une cohorte IceTag-SLS
  contemporaine des mesures capteurs.
\end{enumerate}

Cette démarche repose sur deux hypothèses de travail. L'hypothèse principale est qu'une baisse graduelle touchant plusieurs familles de signaux
est mieux représentée par une détection directionnelle, individualisée et temporelle que par
une succession d'anomalies ponctuelles. L'hypothèse exploratoire est qu'une instabilité
définie par des rapports de mouvement, des transitions et une fragmentation posturale peut
compléter HYPO, à condition de conserver une sortie de surveillance distincte et une
persistance plus courte. Ces hypothèses portent sur le comportement du logiciel et ne
préjugent pas de la cause clinique des événements.

Le système produit ainsi des «~alertes comportementales à vérifier~»: elles soutiennent le suivi
d'élevage sans constituer un diagnostic. Faute d'étalon clinique synchronisé pour toutes les
vaches, le corpus principal sert à étudier rigoureusement la réponse du logiciel à des
perturbations contrôlées et à produire des priorités traçables, dont la portée reste distincte de
celle d'un outil diagnostique validé \citep{alsaaod2019automatic}.

Le corpus principal contient 37\,839 mesures de 28 vaches. Les premiers 60\,\% des intervalles
admissibles de chaque série forment la période de référence, les 40\,\% suivants la période
surveillée. HYPO combine des totaux glissants sur 12 heures, une persistance de six heures, un
score multivarié et une somme cumulative unilatérale (CUSUM). La campagne principale répartit
33 dégradations graduelles et 11 contrôles brefs entre 11 vaches. Chaque perturbation est ajoutée
à une copie de la période future, puis comparée à l'exécution propre de la même vache afin de ne
créditer que les alertes qui lui sont attribuables; les données originales restent intactes.

L'évaluation prend l'événement et la vache comme unités, non chaque intervalle de capteur. Chaque
dégradation forme un épisode dont la détection et la localisation sont mesurées par recouvrement;
les indicateurs sont agrégés par vache et leurs intervalles de confiance proviennent d'un
rééchantillonnage des vaches. Ce choix tient compte de la corrélation entre intervalles successifs.
La campagne compare aussi HYPO aux méthodes Isolation Forest, LOF et pédométrique, relie la
détection à la charge de vérification et scelle les résultats par un manifeste cryptographique.
Un test de stress gelé ajoute cinq formes moins favorables et un contrôle pas-seuls à trois positions
par vache. Leurs aires numériques d'enveloppe sont appariées à celle du profil graduel modéré afin
de contrôler l'intensité, sans supposer une modification physique identique entre vaches ou
familles de signaux.

L'extension INSTABILITÉ agrège trois heures avec une persistance de deux heures. Elle est évaluée
avec HYPO sur 132 scénarios d'intensité, de séquence, de contrôle et de confusion. Une fusion
hiérarchique conserve l'instabilité isolée comme surveillance et réserve la notification à HYPO
ou à une convergence. Cette extension exploratoire enrichit la lecture opérationnelle sans
remplacer l'évaluation technique plus ciblée de HYPO.

Le mémoire apporte ainsi une chaîne reproductible allant du signal brut à une priorité de
vérification. Ses contributions sont les suivantes:
\begin{enumerate}[itemsep=0pt, parsep=0pt, topsep=0.3\baselineskip]
  \item un détecteur HYPO directionnel, persistant, interprétable et individualisé;
  \item un banc d'essai contrefactuel reproductible associant injections post-référence,
  exécution propre et attribution événementielle;
  \item un test de stress à aire d'enveloppe appariée qui distingue la robustesse temporelle de l'apport
  propre de la cohérence multivariée;
  \item une évaluation attribuable par événement et par vache, avec comparaisons à Isolation
  Forest et LOF, ablation pédométrique et analyse OFAT (de l'anglais \emph{One Factor At a
  Time}, soit un facteur à la fois) de 20 variantes;
  \item une extension bidirectionnelle exploratoire avec persistances séparées et cinq
  règles de fusion;
  \item des contrôles brefs et des facteurs confondants qui rendent visibles les limites de spécificité
  au lieu de les masquer;
  \item l'évaluation du temps d'exécution, l'analyse SLS synchronisée, l'interface Streamlit et
  une suite de tests reproductibles.
\end{enumerate}
La contribution principale réside dans l'articulation contrôlée de quatre exigences souvent
traitées séparément: individualisation, temporalité, attribution et reproductibilité. Les
artefacts conservés permettent en outre de reproduire les principaux résultats.

Les six chapitres présentent successivement le contexte de la boiterie
(Chapitre~\ref{chapitre:contexte}), l'état de l'art de sa détection automatisée
(Chapitre~\ref{chapitre:etatart}), le système et son protocole d'évaluation
(Chapitre~\ref{chapitre:systeme}), les corpus et le cadre expérimental
(Chapitre~\ref{chapitre:donnees}), l'interface de priorisation
(Chapitre~\ref{chapitre:interface}), puis les résultats et analyses IceTag-SLS
(Chapitre~\ref{chapitre:resultats}). La discussion et la conclusion distinguent enfin les
résultats établis des observations exploratoires et identifient les travaux nécessaires.

\end{introduction}
